\subsection*{5.2 Independent Random Variables of Discrete Type or Continuous Type}

If $X$ and $Y$ are both discrete random variables, their statistical independence can be expressed in terms of their probability mass functions (pmf's). For simplicity of notation, suppose both random variables $X$ and $Y$ take positive integers as their values. In general, the following theorem holds for any countable codomain.

\begin{thmbox}{5.3}
Suppose $X$ and $Y$ are discrete random variables with pmf's
\[
p_X(m) \triangleq P(X=m) \text{ and } p_Y(n) \triangleq P(Y=n),
\]
respectively, and joint pmf
\[
p_{XY}(m,n) \triangleq P(X=m,\, Y=n)
\]
for $m,n \in \mathbb{N}$. Then $X$ and $Y$ are independent if and only if
\begin{equation}
p_{XY}(m,n)=p_X(m)p_Y(n)
\tag{5.1}
\end{equation}
for all $m,n \in \mathbb{N}$.
\end{thmbox}

\textit{Proof} \\
$(\Rightarrow)$ Let $m$ and $n$ be any positive integers. We apply the definition of independence of random variables with $B_1=\{m\}$ and $B_2=\{n\}$. We obtain (5.1) directly from the definition of independence of random variables.

$(\Leftarrow)$ Conversely, consider two subsets $B_1$ and $B_2$ of $\mathbb{N}$. We have
\[
P(X \in B_1,\, Y \in B_2)=\sum_{m \in B_1}\sum_{n \in B_2} p_{XY}(m,n).
\]

If (5.1) holds, we can compute
\[
P(X \in B_1,\, Y \in B_2)=\sum_{m \in B_1}\sum_{n \in B_2} p_X(m)p_Y(n)
\]
\[
\phantom{P(X \in B_1,\, Y \in B_2)}=\sum_{m \in B_1} p_X(m)\sum_{n \in B_2} p_Y(n)
\]
\[
\phantom{P(X \in B_1,\, Y \in B_2)}=P(X \in B_1)P(Y \in B_2).
\]

This proves that $X$ and $Y$ are independent. \hfill $\square$

For real-valued random variables, we have the following analogous theorem.

\begin{thmbox}{5.4}
Suppose $X$ and $Y$ are random variables on probability space $(\Omega,\mathcal{F},P)$. Let
\[
F_{XY}(x,y) \triangleq P(X \le x,\, Y \le y)
\]
\[
F_X(x) \triangleq P(X \le x)
\]
\[
F_Y(y) \triangleq P(Y \le y).
\]

Then $X$ and $Y$ are independent if and only if
\[
F_{XY}(x,y)=F_X(x)F_Y(y) \quad \text{for all } x,y \in \mathbb{R}.
\]
\end{thmbox}

Like the previous theorem, the forward part is straightforward. The reverse direction requires the $\pi$--$\lambda$ theorem. We establish a convenient result in the next theorem and apply it to prove Theorem 5.4.

\begin{thmbox}{5.5}
Let $(\Omega,\mathcal{F},P)$ be a probability space and $X$ and $Y$ be real-valued random variables defined on $(\Omega,\mathcal{F},P)$. Suppose the Borel algebra $\mathcal{B}(\mathbb{R})$ is generated by a $\pi$-system $\mathcal{C}$. Then $X$ and $Y$ are independent if for all $A,B \in \mathcal{C}$,
\[
P(X \in A,\, Y \in B)=P(X \in A)P(Y \in B).
\]
\end{thmbox}

\textit{Proof} Assume $P(X \in A,\, Y \in B)=P(X \in A)P(Y \in B)$ holds for all $A,B \in \mathcal{C}$.

For any $C \in \mathcal{C}$, define
\[
\mathcal{L}_C \triangleq \{A \in \mathcal{B}(\mathbb{R}) : P(X \in A,\, Y \in C)=P(X \in A)P(Y \in C)\}.
\]

We have $\mathcal{C} \subseteq \mathcal{L}_C$ by assumption, and we want to show that $\mathcal{L}_C$ is a $\lambda$-system. (See Def. 3.10 for the definition of $\lambda$-system.)

(i) Since $P(X \in \Omega,\, Y \in C)=P(Y \in C)=P(X \in \Omega)P(Y \in C)$, we have $\Omega \in \mathcal{L}_C$.

(ii) Suppose $A$ is in $\mathcal{L}_C$. Then,
\[
P(X \in A^c,\, Y \in C)
\]
\[
= P(Y \in C)-P(X \in A,\, Y \in C)
\]
\[
= P(Y \in C)-P(X \in A)P(Y \in C) \qquad \text{(because } A \in \mathcal{L}_C\text{)}
\]
\[
= P(X \in A^c)P(Y \in C).
\]

Therefore $A^c \in \mathcal{L}_C$.

(iii) Suppose $A_i$ is a sequence of mutually disjoint sets in $\mathcal{L}_C$ for $i=1,2,3,\ldots$. Then,
\[
P(X \in \uplus_i A_i,\, Y \in C)=\sum_{i=1}^{\infty} P(X \in A_i,\, Y \in C).
\]

Because $A_i \in \mathcal{L}_C$, we can further simplify it to
\[
P(X \in \uplus_i A_i,\, Y \in C)=\sum_{i=1}^{\infty} P(X \in A_i)P(Y \in C)
\]
\[
= P(Y \in C)P(X \in \uplus_i A_i).
\]

Therefore $\uplus_i A_i \in \mathcal{L}_C$. This proves that $\mathcal{L}_C$ is a $\lambda$-system for each $C \in \mathcal{C}$.

Since $\mathcal{C} \subseteq \mathcal{L}_C$ by assumption and $\mathcal{C}$ generates $\mathcal{B}(\mathbb{R})$, by applying the $\pi$--$\lambda$ theorem, we obtain
\[
\mathcal{B}(\mathbb{R})=\sigma(\mathcal{C}) \subseteq \mathcal{L}_C.
\]

As a result, we have
\begin{equation}
P(X \in B,\, Y \in C)=P(X \in B)P(Y \in C)
\tag{5.2}
\end{equation}
for any fixed $C \in \mathcal{C}$ and all Borel sets $B \in \mathcal{B}(\mathbb{R})$.

Now we fix a Borel set $B \in \mathcal{B}(\mathbb{R})$ and define
\[
\mathcal{L}_B \triangleq \{C \in \mathcal{B}(\mathbb{R}) : P(X \in B,\, Y \in C)=P(X \in B)P(Y \in C)\}.
\]

From what we have proved in (5.2), we have $\mathcal{C} \subseteq \mathcal{L}_B$. By repeating the same steps as in the first part of the proof, we can show that $\mathcal{L}_B$ is a $\lambda$-system for any Borel set $B$. By the $\pi$--$\lambda$ theorem, we have $\mathcal{B}(\mathbb{R}) \subseteq \mathcal{L}_B$ for all $B \in \mathcal{B}(\mathbb{R})$. Hence
\[
P(X \in B,\, Y \in C)=P(X \in B)P(Y \in C)
\]
for all Borel sets $B$ and $C$. \hfill $\square$

We prove Theorem 5.4 by applying the previous theorem with
\[
\mathcal{C}=\{(-\infty,a] : a \in \mathbb{R}\},
\]
which is a $\pi$-system that generates the Borel algebra.

Theorems 5.3 and 5.4 show that a more general definition of independence is backward compatible. The definitions of independence commonly taught in a first course of probability are theorems that can be derived from the first principles.

\textbf{Example 5.2.1 (Independent Gaussian Random Variables)} \\
It is well known that two zero-mean Gaussian random variables, $X_1$ and $X_2$, with zero correlation are independent. Without loss of generality, let us assume that $X_1$ and $X_2$ have zero mean. By ``zero correlation,'' we mean that $E[X_1X_2]=0$. If we set $\rho=0$ in the joint Gaussian probability density function in (1.3), we obtain
\[
f_X(x_1,x_2)=\frac{1}{2\pi\sigma_1\sigma_2}\exp\left(-\frac{x_1^2}{2\sigma_1^2}-\frac{x_2^2}{2\sigma_2^2}\right),
\]
which can be factorized as the product of two Gaussian probability density functions
\[
f_X(x_1,x_2)=\frac{1}{\sqrt{2\pi}\sigma_1}e^{-x_1^2/(2\sigma_1^2)}
\frac{1}{\sqrt{2\pi}\sigma_2}e^{-x_2^2/(2\sigma_2^2)}.
\]

By Theorem 5.4, we conclude that $X_1$ and $X_2$ are independent, which means that for any Borel sets $B_1$ and $B_2$ in $\mathbb{R}$, we have
\[
P(X_1 \in B_1,\, X_2 \in B_2)=P(X_1 \in B_1)\cdot P(X_2 \in B_2).
\]

While uncorrelatedness is necessary for independence, it is a much weaker condition. When the joint probability distribution is not Gaussian, two random variables can be uncorrelated while still being independent.

It is important to note that the term ``jointly'' is also key here. It is possible to construct two dependent and uncorrelated random variables that are Gaussian but not jointly Gaussian. For example, consider two independent random variables: $X \in N(0,1)$, a standard Gaussian random variable, and $U$, a discrete random variable that takes the values $1$ and $-1$ with equal probability. We define a new random variable $Y$ as the product of $U$ and $X$, i.e., $Y=UX$. It follows that $Y$ also follows a standard Gaussian distribution. However, the correlation between $X$ and $Y$ is zero. This illustrates that if two Gaussian random variables have zero mean and zero correlation but are not jointly Gaussian, then they could be dependent.

\textbf{Example 5.2.2 (Function of Dependent Random Variables Can Be Independent)} \\
The converse of Theorem 5.2 is false. We can construct dependent random variables whose squares are independent. Consider the example of jointly distributed random variables $X$ and $Y$, defined by the joint probability density function
\[
f_{XY}(x,y)=\frac{1}{4}(1+xy)
\]
for $(x,y) \in [-1,1] \times [-1,1]$. One can check that the marginal pdf's $f_X(x)$ and $f_Y(y)$ are both constant $1/2$ over the range of $X$ and the range of $Y$, respectively. Hence, $X$ and $Y$ are not independent. However, we can show that
\[
P(X^2 \le x,\, Y^2 \le y)=\sqrt{x}\sqrt{y}=P(X^2 \le x)\cdot P(Y^2 \le y)
\]
for $-1 \le x,y \le 1$. By Theorem 5.4, the random variables $X^2$ and $Y^2$ are independent.
